\par\Needspace{40\baselineskip}\pdfbookmark[1]{Gauss an Schumacher. 1822 November 25.}{schumacher1822}\begin{center}\rule{0.72\linewidth}{0.3pt}\par Astronomische Nachrichten, herausg. v. \textsc{Schumacher}. Nr. 24. Beilage 2. 1822 Dec.\par\rule{0.72\linewidth}{0.3pt}\end{center}

\begin{center}
\textsc{Gauss} an \textsc{Schumacher}.
\end{center}
\begin{flushright}
Göttingen 1822. Nov. 25.
\end{flushright}
--- --- Den von Herrn Professor \textsc{Encke} geäußerten Wunsch nach Logarithmentafeln mit 6 Ziffern habe ich schon öfter selbst gehegt, und Sie werden sich gewiss vielfachen Dank erwerben, wenn Sie solche Tafeln veranlassen. Alles, was Herr \textsc{Encke} über ihre äußere Einrichtung und ihren Inhalt sagt, unterschreibe ich als meine eigene Meinung; nur die Proportionalteile scheinen mir überflüssig, und alles Überflüssige beeinträchtigt den leichten, übersichtlichen Gebrauch. Bei solchen Dingen hängt freilich manches von individueller Gewöhnung ab. Wenn aber einige Benutzer von Tafeln anders gewöhnt sind als ich, so gibt es wohl auch andere mit denselben Gewohnheiten. Daher berühre ich noch einen Umstand, nämlich den Wechsel der \(4^{\mathrm{ten}}\) Ziffer in den Logarithmentafeln. Sie kennen die Einrichtung, die dafür in \textsc{Callet}'s Tafeln getroffen ist und die einige als Verbesserung betrachtet haben. Ich gestehe, dass ich der entgegengesetzten Meinung bin: Die \textit{regelmäßige} Einteilung durch horizontale Linien von 5 zu 5 Zeilen, wie sie in \textsc{Sherwin}'s und anderen Tafeln besteht, halte ich bei häufigem Gebrauch für \textit{viel} wesentlicher und bequemer. Daher mag ich mich auch niemals der \textsc{Callet}'schen Logarithmen bedienen. In meiner langjährigen Praxis kenne ich keinen einzigen Fall, in dem mich der Gebrauch der \textsc{Sherwin}'schen Tafeln bei der \(4^{\mathrm{ten}}\) Ziffer zu einem Rechenfehler verleitet hätte. Daher lege ich auch auf die Sternchen bei \textsc{Vega} und in anderen Tafeln gar keinen Wert und halte mich wegen des besseren Papiers und der schöneren Ziffern lieber an die \textsc{Sherwin}'schen. Wer solche Warnzeichen schätzt, kann in seinem Exemplar an den betreffenden Stellen leicht rote oder grüne Punkte anbringen. ---
